\section{The Control Plane API (FastAPI)}

\subsection{Routers and two kinds of caller}

The API has two audiences, split into separate router files:

\begin{description}
\item[Portal routes] (\texttt{routers/portal.py}) for people: submit jobs,
  list them, stream logs, admin actions. Authenticated by a session cookie from
  the university login, or the \texttt{X-Dev-User} header in development.
\item[Agent routes] (\texttt{routers/agent.py}) for machines: register,
  heartbeat, claim, renew, logs, outputs, complete. Authenticated by a
  per-node bearer token.
\end{description}

\subsection{Dependency injection}

FastAPI lets an endpoint declare what it needs as parameters, and builds those
values before calling it. Authentication is written once as a dependency:

\begin{lstlisting}[language=Python]
# control/app/deps.py (simplified)
async def get_current_node(
    session: AsyncSession = Depends(get_session),
    authorization: str | None = Header(default=None),
    x_boot_session_id: str | None = Header(default=None,
                                           alias="X-Boot-Session-Id"),
) -> Node:
    token = authorization.split(" ", 1)[1]
    cred = await session.scalar(select(NodeCredential).where(
        NodeCredential.token_hash == hash_token(token),
        NodeCredential.revoked_at.is_(None)))
    if cred is None:
        raise AppError(401, "auth", "invalid node token")
    node = await session.get(Node, cred.node_id)
    if x_boot_session_id != node.boot_session_id:
        raise AppError(409, "stale_boot_session", "re-register")
    return node

@router.post("/agent/heartbeat")
async def heartbeat(body: HeartbeatIn,
                    node: Node = Depends(get_current_node), ...):
    ...  # only runs for an authenticated, current node
\end{lstlisting}

Every agent endpoint gets authentication, revocation and boot-session checks
just by listing \texttt{node: Node = Depends(get\_current\_node)}.

\paragraph{Request validation.} Request bodies are Pydantic models
(\texttt{schemas.py}). A heartbeat with a missing field or a string where a
number belongs is rejected with \texttt{422} before any code runs.

\subsection{Long polling in code}

\begin{lstlisting}[language=Python]
@router.post("/agent/claim")
async def claim_job(node: Node = Depends(get_current_node),
                    session: AsyncSession = Depends(get_session)):
    waited = 0.0
    while True:
        attempt = await claim(session, node)
        if attempt:
            payload = await claim_payload(session, attempt)
            await session.commit()
            return {"attempt": payload, "build": None}
        await session.commit()          # release any locks while waiting
        if waited >= settings.claim_poll_seconds:
            return {"attempt": None, "build": None}
        await asyncio.sleep(1.0)        # costs almost nothing in async code
        waited += 1.0
\end{lstlisting}

The \texttt{commit()} before sleeping matters: holding a transaction open while
asleep would hold locks and a database connection for no reason.

\subsection{Live logs with Server-Sent Events}

The student's browser shows log lines as they arrive. The portal uses
\textbf{Server-Sent Events} (SSE): an ordinary HTTP response that never ends,
where the server writes one small text block per event:

\begin{verbatim}
id: 42
event: log
data: {"seq": 42, "stream": "stdout", "line": "epoch 3 loss 0.41"}

id: 43
event: log
data: {"seq": 43, "stream": "stdout", "line": "epoch 4 loss 0.37"}
\end{verbatim}

If the connection drops, the browser reconnects automatically and sends
\texttt{Last-Event-ID: 43}; the server resumes after line 43. That only works
because every log line has a sequence number, assigned by the agent and checked
by the server (a chunk must start at \texttt{last\_log\_seq + 1}).

\paragraph{SSE or WebSockets?} SSE is one-way (server to browser), plain HTTP,
and reconnects by itself, which is all a log viewer needs. WebSockets are
two-way, useful for chat or collaborative editing. Other SSE users: ChatGPT-style
token streaming, stock tickers, build logs in CI dashboards.

\subsection{Background workers}

Two loops start with the app (\texttt{main.py}):

\begin{itemize}
\item The \textbf{lease worker}, every 2 seconds: expire leases and mark silent
  nodes offline.
\item The \textbf{alert worker}, every 60 seconds: compute alerts and post new
  ones to a webhook.
\end{itemize}

They are \texttt{asyncio} tasks in the same process, created in FastAPI's
\texttt{lifespan} hook and cancelled on shutdown. For one server that is enough;
a larger system would move them to a separate worker process or a scheduler
like Celery beat.

\subsection{Errors as data}

All expected failures raise \texttt{AppError(status, code, message)}, which a
single handler turns into JSON like
\texttt{\{"error": "fenced", "message": "stale fencing token"\}}. Clients
branch on the stable \texttt{error} code, never on the human message. The agent,
for example, re-registers when it sees \texttt{stale\_boot\_session}.
